% Round 12. Explicit scope clarification; no historical source is edited.
\section{The supported-complex construction over a nongeometrically integral base}
\label{r12i:section}

\begin{lemma}[Native integral scope of the supported-complex construction]
\label{r12i:integral-PC}
Keep the original finite extension $K/\mathbf Q_p$ and its absolute
value.  The Round 11 supported-complex theorem, the exact-support
Green-pair construction for every prescribed rational cycle, and
the generic-Chow-zero global primitive construction apply to every
smooth integral projective $K$-variety $X$.  Geometric integrality
over $K$ is not required.  Every actual form, test and finite-trace
operation in these statements is in the original $K$-analytic theory.
\end{lemma}

\begin{proof}
The supported-complex section and its main theorem already assume
only that $X$ is smooth and integral.  We check the one potentially
sensitive ingredient and then the passage to an arbitrary prescribed
cycle, whose Round 11 corollary was stated with the stronger hypothesis.

Let $F\to Y$ be the smooth projective flag fibration appearing after
the kernel-Grassmann modification and its characteristic-zero generic
resolution.  Both $F$ and $Y$ are smooth integral $K$-varieties.
Put $d=\dim F-\dim Y$.  If $d=0$, the map is already finite etale
and its integral source can be used.  If $d>0$, the uniform fiber
dimension argument in Round 11 gives a nonempty open of tuples of
high-degree very ample sections whose common zero locus has finite
nonempty fibers over every point of $Y$.  Its dimension estimate is
made on every geometric fiber and does not require the total space
or the base to be geometrically integral.

Intersect this parameter open with the smooth complete-intersection
Bertini open.  The latter requires smoothness of the source and a
linear system giving an immersion, not geometric irreducibility;
these are exactly the hypotheses of Stacks Lemma 33.47.3,
Tag~\texttt{0FD6}.  Successive restrictions of the very ample linear
system still give immersions.  Equivalently, after algebraic base
change the smoothness conditions are imposed simultaneously on the
finitely many geometric components; the resulting nonempty open is
defined over $K$.  Since $K$ is infinite it has a $K$-rational tuple.
No analytic scalar extension is involved in this algebraic check.

The resulting complete intersection $T_0^{\rm all}$ is smooth,
projective, and of pure dimension $\dim Y$.  Its map to $Y$ is
proper with finite nonempty fibers, hence finite surjective.
Miracle flatness applies, since its source is Cohen--Macaulay and
$Y$ is regular; see Stacks Tag~\texttt{00R4}.
Every connected component of a regular scheme is integral, because
its irreducible components are disjoint.  Choose any nonempty
connected component $T_0$.  It is open and closed in
$T_0^{\rm all}$, so its finite map is still flat.  Its image is both
open and closed in integral $Y$, and therefore all of $Y$.
Consequently $T_0\to Y$ is smooth-source, integral, finite flat and
surjective, with its own native positive degree.

Choosing this component does not spoil the supported $K_0$ identity:
the identity is pulled back to the chosen component, rather than
being an equality that would require summing the classes of all
components of a fixed intersection cycle.  Subsequent target-flag
graph closures and generic resolutions are birational on this
integral source and preserve its full prescribed finite open.
The leading supported-character identity then uses
$f_*[T]=d[X]$ with this actual function-field degree $d$.
Normality needed by the minimal trace follows from smoothness of
the generic source and target.  None of these statements requires
geometric integrality.

For a prescribed cycle $Z=\sum_j a_j[Z_j]$ on this $X$, resolve each
$\mathcal O_{Z_j}$ by finitely many generic vector bundles.  Such
resolutions use regularity and the projective resolution property,
which hold for smooth integral $X$ over $K$.  Extend the bundles by
the same finite Grassmann-graph construction to a projective flat
model, and extend the differentials after uniformizer rescaling.
Each model complex has generic support exactly $Z_j$.  The preceding
supported-complex theorem applies component by component.  Summing
with the nonzero coefficients gives a Green pair on exactly
$X^{\rm an}\setminus|Z|^{\rm an}$, with the original cycle and all
actual tests retained.  This proves the stated extension of the
exact-support corollary.

The global primitive proof uses the same construction with generic
open $U=X$.  Generic $K_0$ relations, Grassmann subbundle graphs,
and rational Adams projectors require smooth integral projectivity,
not geometric integrality.  Their primitives consequently retain
the original global scalar-curvature line sums and finite-flat maps
over $K$, exactly as in Round 11.
\end{proof}

\begin{remark}[Finite constant extensions are treated as native schemes]
Let $K'/K$ be finite and let $Y/K'$ be smooth integral projective.
The composite $Y\to\operatorname{Spec}K$ is smooth because
$K'/K$ is separable, and it is still integral and projective.
A projective flat $R'$-model of $Y$ is projective and flat over $R$.
Lemma~\ref{r12i:integral-PC} can therefore be run directly on this
$R$-model, producing $K$-actual forms and testing against arbitrary
$K$-actual tests.  This is not an inference from a formula tested
only over $K'$.

The special fiber over $R$ keeps its ramification thickening, and
the degree of a finite generic map is its full degree as a
$K$-morphism.  Thus no comparison of native skeleton measures over
different ground fields is required for this scope clarification.
The same native treatment was already stated in the proof of
Round 11's finite good-cover corollary; the present lemma makes its
extension to the exact-support corollary explicit.
\end{remark}
